% 修改对应部分就可以进行整个文档编辑

% 注意：下面每个 \newcommand 定义的都是“字符串内容”，改等号后面花括号里的东西即可，不需要则留空

% 成绩
\newcommand{\infoScore}{}




%=========================================================================
% 
%                       课    程    信    息
%
%=========================================================================
% 实验题目
\newcommand{\infoExpNo}{实验N}
\newcommand{\infoExpTitle}{实验名}
% 课程名称
\newcommand{\infoCourse}{课程名}
% 实验日期（格式 2026-09-01）
\newcommand{\infoExperimentTime}{实验日期}
% 任课教师姓名
\newcommand{\infoTeacher}{教师名}
% 总学时
\newcommand{\infoHours}{学时}
% 小组编号
\newcommand{\infoGroup}{编号}




%=========================================================================
% 
%                       个    人    信    息
%
%=========================================================================

% 学生姓名
\newcommand{\infoName}{李四}
% 学号
\newcommand{\infoStudentId}{20240001}
% 学院名称
\newcommand{\infoCollege}{信息科学技术学院}
% 班级名称
\newcommand{\infoClass}{aaa}
% 班级人数
\newcommand{\infoStudentCount}{60}
% 学年开始年份
\newcommand{\infoYearStart}{26}
% 学年结束年份
\newcommand{\infoYearEnd}{27}
% 学期
\newcommand{\infoSemester}{1}
% 专业名称
\newcommand{\infoMajor}{aaa}
% 年级（24 = 2024 级）
\newcommand{\infoGrade}{24}
% 班级序号（2 = 2 班）
\newcommand{\infoClassNo}{2}

